\documentclass[20pt,margin=1in,innermargin=-4.5in,blockverticalspace=-0.25in]{tikzposter}
\geometry{paperwidth=33.11in,paperheight=46.81in} %A0
% \geometry{paperheight=33.11in,paperwidth=23.4in} %A1
\usepackage[utf8]{inputenc}
\usepackage{csquotes}
%\usepackage[estonian]{babel}
\usepackage{amsmath}
\usepackage{amsfonts}
\usepackage{amsthm}
\usepackage{amssymb}
\usepackage{mathrsfs}
\usepackage{graphicx}
\usepackage{lipsum}
\usepackage[export]{adjustbox}
\usepackage{enumitem}
\usepackage[backend=biber,style=numeric]{biblatex}
\usepackage{unitartu-theme}
\makeatletter
\setlength{\TP@visibletextwidth}{31.0in}
\setlength{ \TP@visibletextheight}{45in}
\makeatother
\usepackage{mwe} % for placeholder images
\usepackage{bm}
\usepackage{bbm}

\usepackage{tikz}


% siunit unit
\usepackage{siunitx}

% comments
\usepackage{verbatim}

% for footnotes in tables
\usepackage{booktabs}
%\usepackage{caption}
\usepackage[flushleft]{threeparttable}


% siunit same text as base text
\sisetup{detect-all}

% degree sign
\usepackage{gensymb}

% angstrom
\newcommand{\angstrom}{\mbox{\normalfont\AA}}

\addbibresource{refs.bib}

% set theme parameters
\tikzposterlatexaffectionproofoff
\usetheme{UniTartuTheme}
\usecolorstyle{UniTartuStyle}

\usepackage[scaled]{helvet}
\renewcommand\familydefault{\sfdefault} 
\renewcommand{\vec}[1]{\bm{#1}}
\newcommand{\Tr}{\text{Tr}}
\usepackage[T1]{fontenc}
\usepackage{tgbonum}


\title{\parbox{0.8\linewidth}{Laser shock peening acoustic emission measurement by evaluation of acoustic wave propagating through air}}
\author{\textbf{Marek Böhm}\textsuperscript{1,2}, \textbf{Sunil Pathak}\textsuperscript{1},\textbf{Jan Kaufman}\textsuperscript{1}, \textbf{Jan Šmaus}\textsuperscript{1}, \textbf{Ondřej Stránský}\textsuperscript{3}}

\institute{\textsuperscript{1}HiLASE Centre, Institute of Physics of the Czech Academy of Sciences, Czech Republic 
\newline
\textsuperscript{2}Czech Technical University in Prague, Faculty of Nuclear Sciences and Physical Engineering, Czech Republic
\newline
\textsuperscript{3}Czech Technical University in Prague, Faculty of Mechanical Engineering, Czech Republic}

\titlegraphic{\includegraphics[width=0.2\linewidth]{Logo_HiLASE_2C_Pantone.pdf}}

% begin document
\begin{document}
\maketitle
\centering
\begin{columns}
    \column{0.5}
\block{Abstract}{
    Acoustic emission (AE) monitoring is a non-destructive evaluation technique widely used to detect and analyse the acoustic waves generated by material deformation and damage. In this study, we applied AE monitoring to laser shock peening (LSP). LSP treatment technique uses a high-intensity laser beam to generate a shock wave that induces compressive residual stresses in metallic materials. We conducted LSP experiments on alloy samples and simultaneously recorded the AE signals generated during the process. The AE signals were analysed to identify the different types of acoustic events and their corresponding sources, including breakdown in water and dielectric breakdown within the laser beam in the air. The results show that AE monitoring can provide valuable information on the LSP process and its effects on the material properties. Furthermore, the AE signals can be used to optimize the LSP parameters, such as the laser energy and spot size, to achieve the desired compressive residual stresses and avoid undesired damage. Overall, this study demonstrates the potential of AE monitoring as a powerful tool for the in-situ characterization of LSP, providing insights into the underlying physical processes and helping to improve the quality and reliability of the treated materials.
     
}
\block{Experimental methods}{

    %%%%%%%%%%%%%%%%%%%%%%%%%%%%%%%%%%%%%%%%%%%%%%%%%%%%%%%%%%%%%%%%%
    %%%%%%%%%%%%%%%%%%%%%%%%%%%%%%%%%%%%%%%%%%%%%%%%%%%%%%%%%%%%%%%%%
    %%%%%%%%%%%%%%%%%%%%%%%%%%%%%%%%%%%%%%%%%%%%%%%%%%%%%%%%%%%%%%%%%
    % MATERIAL

    \section*{Material}
    Aluminium AA2024 was machined into dimensions of \(\SI[]{100}{} \times \SI[]{100}{} \times \SI[]{5}{}\) \SI[]{}{\mm}. The surface roughness of the material was \( R_a  = \SI[]{0.754}{\micro\metre} \). The surface was not polished.  The surface residual stress baseline level measured by X-ray diffraction was  \( \SI[]{}{\sigma}  = \SI[separate-uncertainty = true]{8.15 \pm 7.80 }{\mega\pascal} \).


    %%%%%%%%%%%%%%%%%%%%%%%%%%%%%%%%%%%%%%%%%%%%%%%%%%%%%%%%%%%%%%%%%
    %%%%%%%%%%%%%%%%%%%%%%%%%%%%%%%%%%%%%%%%%%%%%%%%%%%%%%%%%%%%%%%%%
    %%%%%%%%%%%%%%%%%%%%%%%%%%%%%%%%%%%%%%%%%%%%%%%%%%%%%%%%%%%%%%%%%
    % LASER SYSTEM

    \section*{Laser system}
    A Litron LPY ST 7875-10 2HG laser was used for peening the samples. The Litron laser is a pulsed Q-switched Nd:YAG laser with a oscillator-amplifier configuration with stable resonators. The system generates pulses tens of nanoseconds in duration at \SI{532}{\nano\meter} with a pulse energy up to \SI{1.5}{\joule} at \SI{10}{\hertz} repetition rate. The critical parameters of this system are summarized in Table~1~\cite{litron}. The temporal shape of the Litron laser pulse measured by an InGaAs EOT ET-3000EXT PIN detector, recorded by a Rohde\&Schwarz RTA4004 Digital Oscilloscope and the profile of the laser pulse at \SI{170}{\milli\joule} recorded by an Allied Vision Manta G-210B ASG GigE camera are shown in Figure \ref{fig:profile}. 
    The FWHM of the laser pulse is \SI{18.6}{\ns}.


    %%%%%%%%%%%%%%%%%%%%%%%%%%%%%%%%%%%%%%%%%%%%%%%%%%%%%%%%%%%%%%%%%
    %%%%%%%%%%%%%%%%%%%%%%%%%%%%%%%%%%%%%%%%%%%%%%%%%%%%%%%%%%%%%%%%%
    %%%%%%%%%%%%%%%%%%%%%%%%%%%%%%%%%%%%%%%%%%%%%%%%%%%%%%%%%%%%%%%%%
    % LASER PROFILE FIGURE

        % width for 
        % you can add captions by tikzfigure
        \begin{center}
        \begin{tikzfigure}[Temporal shape of Litron system laser pulse] 
        \begin{minipage}{.4\textwidth}
        \includegraphics[width=0.45\linewidth,valign=c]{img/Litron_temporal_plotly.png}
        \label{fig:profile}
        \end{minipage}%
        \begin{minipage}{.4\textwidth}
        \includegraphics[width=0.45\linewidth,valign=c]{img/C2-v3_36_mJ_scale.png}
        \label{fig:profile}
        \end{minipage}%
        \end{tikzfigure}
        \end{center}


    %%%%%%%%%%%%%%%%%%%%%%%%%%%%%%%%%%%%%%%%%%%%%%%%%%%%%%%%%%%%%%%%%
    %%%%%%%%%%%%%%%%%%%%%%%%%%%%%%%%%%%%%%%%%%%%%%%%%%%%%%%%%%%%%%%%%
    %%%%%%%%%%%%%%%%%%%%%%%%%%%%%%%%%%%%%%%%%%%%%%%%%%%%%%%%%%%%%%%%%
    % LASER PARAMETERS
 

    \vspace{1em}

    \begin{center}  

    \begin{threeparttable}
        \centering
        \begin{tabular}{|c | c|} 
        \hline
            \textbf{Parameter} & \textbf{Value} \\ [0.5ex] 
        \hline
        Repetition Rate [Hz] & 10  \\ 
        \hline
            Output Energy [mJ] & \\
            1064 nm & 3500  \\
            532 nm & 1750 \\
        \hline
            Beam Diameter [mm] & 15  \\
        \hline
            Beam Divergence [mrad] & \textless 0.5  \\ 
        \hline
            Pulse Width @1064 nm [ns] & 10-12 \\
        \hline
            Pointing Stability [µrad] & 25  \\
        \hline
            Timing jitter [ns] & \textless 0.5  \\
        \hline
        \end{tabular}

        \caption[Litron~LPY~ST~7875-10~2HG parameters]{Litron~LPY~ST~7875-10~2HG parameters \protect\cite{litronmanual}}
        
       
    \end{threeparttable}

    \label{tab:litronparameters}
    \end{center}

    \section*{Optical setup}
    filler text


    %%%%%%%%%%%%%%%%%%%%%%%%%%%%%%%%%%%%%%%%%%%%%%%%%%%%%%%%%%%%%%%%%
    %%%%%%%%%%%%%%%%%%%%%%%%%%%%%%%%%%%%%%%%%%%%%%%%%%%%%%%%%%%%%%%%%
    %%%%%%%%%%%%%%%%%%%%%%%%%%%%%%%%%%%%%%%%%%%%%%%%%%%%%%%%%%%%%%%%%
    % OPTICAL SETUP SCHEMATICS

        % width for 
        % you can add captions by tikzfigure
        \begin{tikzfigure}[Laser shock peening optical setup schematics]        
        \includegraphics[width=0.5\textwidth,valign=c]{img/optical setup final_v_7.png}
        \label{fig:mesh1}
        \end{tikzfigure}








}

    %%%%%%%%%%%%%%%%%%%%%%%%%%%%%%%%%%%%%%%%%%%%%%%%%%%%%%%%%%%%%%%%%
    %%%%%%%%%%%%%%%%%%%%%%%%%%%%%%%%%%%%%%%%%%%%%%%%%%%%%%%%%%%%%%%%%
    %%%%%%%%%%%%%%%%%%%%%%%%%%%%%%%%%%%%%%%%%%%%%%%%%%%%%%%%%%%%%%%%%
    %%%%%%%%%%%%%%%%%%%%%%%%%%%%%%%%%%%%%%%%%%%%%%%%%%%%%%%%%%%%%%%%%
    %%%%%%%%%%%%%%%%%%%%%%%%%%%%%%%%%%%%%%%%%%%%%%%%%%%%%%%%%%%%%%%%%
    %%%%%%%%%%%%%%%%%%%%%%%%%%%%%%%%%%%%%%%%%%%%%%%%%%%%%%%%%%%%%%%%%
    % SECOND COLUMN




% other side of the page - 2nd column
    \column{0.5}

    \block{Experimental methods (cont.)}{

    %%%%%%%%%%%%%%%%%%%%%%%%%%%%%%%%%%%%%%%%%%%%%%%%%%%%%%%%%%%%%%%%%
    %%%%%%%%%%%%%%%%%%%%%%%%%%%%%%%%%%%%%%%%%%%%%%%%%%%%%%%%%%%%%%%%%
    %%%%%%%%%%%%%%%%%%%%%%%%%%%%%%%%%%%%%%%%%%%%%%%%%%%%%%%%%%%%%%%%%
    % RESIDUAL STRESS MEASUREMENT

    \section*{Residual stress measurement}
    The residual stress measurements were carried out using the hole drilling technique  and X-ray diffraction using conventional X-ray diffraction \(\sin ^2(\Psi )\) technique. The hole drilling system PRISM by Stresstech company was used to gain depth-resolved residual stress profiles. The hole drilling method is based on the measurements of surface distortion caused by the drilling using electronic speckle pattern interferometry. The drill diameter was 3.2 mm, and it enabled measurements to 1.9 mm in depth.  A Rigaku XRD Automate 2 instrument and electrolytic removal was used for analyzing residual stresses in two orthogonal directions. XRD residual stress measurement parameters are provided in Table~2. 

    %%%%%%%%%%%%%%%%%%%%%%%%%%%%%%%%%%%%%%%%%%%%%%%%%%%%%%%%%%%%%%%%%
    %%%%%%%%%%%%%%%%%%%%%%%%%%%%%%%%%%%%%%%%%%%%%%%%%%%%%%%%%%%%%%%%%
    %%%%%%%%%%%%%%%%%%%%%%%%%%%%%%%%%%%%%%%%%%%%%%%%%%%%%%%%%%%%%%%%%
    % XRD PARAMETERS
    \begin{center}  

    \begin{threeparttable}
        \centering
        \begin{tabular}{|c | c|} 
        \hline
            \textbf{Item} & \textbf{Description} \\ [0.5ex] 
        \hline
            Radiation & CrK\(\alpha_1\) \SI{}{\lambda} = \SI{2.2897}{\angstrom} \\
        \hline
            Tilt Angles [\degree] &  0, \SI{\pm2.58}{}, \SI{\pm9.07}{}, \SI{\pm12.45}{}, \SI{\pm18.8}{}, \SI{\pm23.0}{}  \\ 
        \hline
            Aperture size(diameter) [mm] & 2 \\
        \hline
            Bragg's Angle [\degree]  & 139\degree \\

        \hline
        \end{tabular}

        \caption[Litron~LPY~ST~7875-10~2HG parameters]{Parameters for XRD residual stress measurement}
        
       
    \end{threeparttable}

    \label{tab:xrdparameters}
    \end{center}


    %%%%%%%%%%%%%%%%%%%%%%%%%%%%%%%%%%%%%%%%%%%%%%%%%%%%%%%%%%%%%%%%%
    %%%%%%%%%%%%%%%%%%%%%%%%%%%%%%%%%%%%%%%%%%%%%%%%%%%%%%%%%%%%%%%%%
    %%%%%%%%%%%%%%%%%%%%%%%%%%%%%%%%%%%%%%%%%%%%%%%%%%%%%%%%%%%%%%%%%
    % LSP processing
    
    \section*{LSP processing}

    The samples were LSP processed with parameters listed in Table~3. 

    %%%%%%%%%%%%%%%%%%%%%%%%%%%%%%%%%%%%%%%%%%%%%%%%%%%%%%%%%%%%%%%%%
    %%%%%%%%%%%%%%%%%%%%%%%%%%%%%%%%%%%%%%%%%%%%%%%%%%%%%%%%%%%%%%%%%
    % LSP processing parameters
    \begin{center}  

    \begin{threeparttable}
        \centering
        \begin{tabular}{|c|c | c| c| c| c| c|} 
        \hline
            \textbf{Distance [\SI{}{\m}]} & \textbf{E [\SI{}{\milli\joule}]} & \textbf{t [\SI{}{\ns}]} & \textbf{Dimple Ø [\SI{}{ \mm}]}  & \textbf{Overlap [\%] } & \textbf{No. of layers [-] } & \textbf{\( I_0  [\SI[]{}{\giga\watt/\cm^2} \)] }\\ [0.5ex] 
  
        \hline
          0.756 &  800   & 18.6 & 0.893 & 50 & 1 & 6.87  \\

        \hline
        \end{tabular}

        \caption[Litron~LPY~ST~7875-10~2HG parameters]{Parameters for XRD residual stress measurement}
        
       
    \end{threeparttable}

    \label{tab:xrdparameters}
    \end{center}

    }

    \block{Conclusions}{
        \lipsum[100]
        \cite{dirac} \cite{knuthwebsite}
        \lipsum[100]
        \lipsum[100]
    }

    \block{Acknowledgements}{
        This work was supported by European Union’s Horizon 2020 research and innovation programme [Grant No. 739573]. 
    }
        \block{References}{
                \begin{center}
                   \mbox{}\vspace{-1\baselineskip}
    \printbibliography[heading=none] 
        \end{center}
        }

\end{columns}
\end{document}